%! Transformations of Functions — Unit 1, Lesson 1.4 — Guided Notes
\documentclass[10pt]{article}
\usepackage{precalculus-article}
\usepackage{precalculus-boxes}
\newcommand{\TallMath}[1]{$\displaystyle #1\rule[-1.4em]{0pt}{3.2em}$}

\begin{document}

\pageheader{Unit 1, Lesson 1.4}{Guided Notes}

\vspace{0.12in}

\begin{objectivebox}
By the end of this lesson, I will be able to\ldots
\begin{enumerate}[itemsep=2pt]
  \item Sort a transformation rule into a change to the \blank{2.4cm} I hand
        the parent function or a change to the \blank{2.4cm} it hands back.
  \item Describe a vertical translation, a vertical dilation, and a
        \blank{2.4cm} from the numbers \textit{outside} the function.
  \item Explain why a change \textit{inside} the parentheses moves the graph
        in the \blank{2.4cm} direction.
  \item Follow a \blank{2.4cm} point from the parent to the transformed graph,
        and read off the new \blank{2cm} and \blank{1.6cm}.
\end{enumerate}
\end{objectivebox}

\vspace{0.1in}

\boxguard
\begin{vocabbox}
Fill in each term as we define it together.
\par\vspace{2pt}
\noindent\textbf{\textcolor{plum}{Parent function:}}\\[1pt]
\writeline\\[3pt]
\noindent\textbf{\textcolor{plum}{Transformation:}}\\[1pt]
\writeline\\[3pt]
\noindent\textbf{\textcolor{plum}{Translation:}}\\[1pt]
\writeline\\[3pt]
\noindent\textbf{\textcolor{plum}{Dilation:}}\\[1pt]
\writeline\\[3pt]
\noindent\textbf{\textcolor{plum}{Reflection:}}\\[1pt]
\writeline
\end{vocabbox}

\vspace{0.1in}

\boxguard[20]
\begin{hookbox}
$D(t)$ is the water depth at North Harbor, in feet, $t$ hours after midnight.
South Harbor rises and falls exactly the same way, and the harbormaster
publishes its depth as
\[
  S(t) = D(t + 3).
\]

\smallskip
\begin{center}
\renewcommand{\arraystretch}{1.4}
\begin{tabular}{|c|c|c|c|c|c|}
\hline
$t$ (hours after midnight) & 0 & 3 & 6 & 9 & 12 \\
\hline
$D(t)$ (feet) & 4 & 8 & 12 & 8 & 4 \\
\hline
\end{tabular}
\end{center}
\smallskip

The rule \textit{adds} three hours. Does South Harbor's high tide come
\textbf{before} or \textbf{after} North Harbor's? Commit to an answer before
you compute anything.
\writelines{2}
\end{hookbox}

\vspace{0.1in}

\boxguard[24]
\begin{notesbox}{1. The Question and the Answer}
Every rule today is built out of one \textbf{parent function}. Ours is
$f(x) = x^2$.

\smallskip
\begin{center}
\renewcommand{\arraystretch}{1.4}
\begin{tabular}{|c|c|c|c|c|c|}
\hline
$x$ & $-2$ & $-1$ & 0 & 1 & 2 \\
\hline
$f(x) = x^2$ & 4 & 1 & 0 & 1 & 4 \\
\hline
\end{tabular}
\end{center}
\smallskip

Using $f$ is a two-step transaction: you hand $f$ an input, and $f$ hands you
back an output. So there are exactly two places a rule can interfere.

\medskip
\begin{center}
\begin{tikzpicture}
  \node[font=\large, anchor=west] (gq) at (0,0) {$g(x) \;=\;$};
  \node[font=\large, goldacc, anchor=west] (a) at (gq.east) {$a$};
  \node[font=\large, anchor=west] (mid) at (a.east) {$\cdot\, f($};
  \node[font=\large, plum, anchor=west] (inner) at (mid.east) {$\,x + h\,$};
  \node[font=\large, anchor=west] (close) at (inner.east) {$)$};
  \node[font=\large, goldacc, anchor=west] (k) at (close.east) {$\;+\;k$};
  \node[font=\scriptsize\bfseries, plum, align=center]
    (ql) at ([yshift=-0.75cm]inner.south) {inside: the question you hand $f$};
  \draw[thick, plum, ->] (ql.north) -- (inner.south);
  \node[font=\scriptsize\bfseries, goldacc, align=center]
    (al) at ([yshift=0.62cm]mid.north) {outside: the answer $f$ hands back};
  \draw[thick, goldacc, ->] (al.west) to[out=180, in=105] (a.north);
  \draw[thick, goldacc, ->] (al.east) to[out=0, in=80] (k.north);
\end{tikzpicture}
\end{center}

\medskip
\begin{itemize}[itemsep=4pt]
  \item A change to the \textbf{answer} is \blank{2.4cm} and does exactly what
        it says.
  \item A change to the \textbf{question} is \blank{2.4cm}, and it runs
        \blank{2.4cm} to what it looks like --- because you have to work out
        what $x$ makes the question come out right.
\end{itemize}

\medskip
Sort each rule before computing anything.

\smallskip
\begin{center}
\renewcommand{\arraystretch}{1.5}
\begin{tabular}{|c|c|c|}
\hline
\textbf{rule} & \textbf{question or answer?} & \textbf{horizontal or vertical?} \\
\hline
$g(x) = f(x) - 5$ & \blank{2.4cm} & \blank{2.4cm} \\
\hline
$g(x) = f(x - 5)$ & \blank{2.4cm} & \blank{2.4cm} \\
\hline
$g(x) = 4f(x)$ & \blank{2.4cm} & \blank{2.4cm} \\
\hline
$g(x) = f(4x)$ & \blank{2.4cm} & \blank{2.4cm} \\
\hline
\end{tabular}
\end{center}
\end{notesbox}

\vspace{0.1in}

\boxguard[20]
\begin{notesbox}{2. Rewriting the Answer --- Vertical Moves}
Whatever $f$ hands back, do something to it. The input column is never
touched.

\smallskip
\begin{center}
\renewcommand{\arraystretch}{1.4}
\begin{tabular}{|c|c|c|c|c|c|}
\hline
$x$ & $-2$ & $-1$ & 0 & 1 & 2 \\
\hline
$f(x) = x^2$ & 4 & 1 & 0 & 1 & 4 \\
\hline
$f(x) + 3$ & \blank{1cm} & \blank{1cm} & \blank{1cm} & \blank{1cm} & \blank{1cm} \\
\hline
$2f(x)$ & \blank{1cm} & \blank{1cm} & \blank{1cm} & \blank{1cm} & \blank{1cm} \\
\hline
$-f(x)$ & \blank{1cm} & \blank{1cm} & \blank{1cm} & \blank{1cm} & \blank{1cm} \\
\hline
\end{tabular}
\end{center}
\smallskip

\begin{center}
\begin{tikzpicture}[x=0.72cm, y=0.28cm]
  \draw[linegray, very thin] (-3,-5) grid (3,8);
  \draw[->, thick] (-3.5,0) -- (3.7,0) node[below right, font=\small] {$x$};
  \draw[->, thick] (0,-5.6) -- (0,8.8) node[above, font=\small] {$y$};
  \foreach \x in {-2,-1,1,2}
    \node[below, font=\tiny, fill=white, inner sep=1pt] at (\x,0) {\x};
  \foreach \y in {-4,4,8}
    \node[left, font=\tiny, fill=white, inner sep=1pt] at (0,\y) {\y};
  \draw[thick, goldacc, dashed, smooth, samples=60, domain=-2.2:2.2]
    plot (\x, {\x*\x + 3});
  \draw[thick, slate, dashed, smooth, samples=60, domain=-2.2:2.2]
    plot (\x, {-\x*\x});
  \draw[very thick, plum, smooth, samples=60, domain=-2.7:2.7]
    plot (\x, {\x*\x});
  \node[font=\scriptsize, goldacc, fill=white, inner sep=1pt] at (-1.55,7.6) {$f(x)+3$};
  \node[font=\scriptsize, plum, fill=white, inner sep=1pt] at (1.75,7.6) {$f(x)$};
  \node[font=\scriptsize, slate, fill=white, inner sep=1pt] at (1.7,-4.3) {$-f(x)$};
\end{tikzpicture}
\end{center}

\begin{itemize}[itemsep=4pt]
  \item $g(x) = f(x) + k$ is a \textbf{vertical translation}: every answer
        gains $k$. If $k > 0$ the graph moves \blank{1.4cm}; if $k < 0$ it
        moves \blank{1.4cm}.
  \item $g(x) = af(x)$ is a \textbf{vertical dilation} by a factor of
        \blank{1.4cm}: every answer is multiplied by $a$. If $|a| > 1$ the
        graph is stretched away from the $x$-axis; if $0 < |a| < 1$ it is
        \blank{2.4cm} toward it.
  \item If $a < 0$ the dilation also carries a \textbf{reflection} over the
        \blank{1.6cm}-axis. The case $a = -1$ is the flip on its own.
  \item The questions were never touched, so the \blank{2cm} does not change.
        The \blank{1.6cm} does.
\end{itemize}
\end{notesbox}

\vspace{0.1in}

\boxguard[20]
\begin{notesbox}{3. Rewriting the Question --- Horizontal Moves}
Now the rule changes what you hand $f$ before $f$ ever goes to work.

\medskip
\textbf{(a) $g(x) = f(x + 3)$.} Fill each cell by first saying what gets handed
to $f$.

\smallskip
\begin{center}
\renewcommand{\arraystretch}{1.4}
\begin{tabular}{|c|c|c|c|c|c|}
\hline
$x$ & $-5$ & $-4$ & $-3$ & $-2$ & $-1$ \\
\hline
handed to $f$: $x + 3$ & \blank{1cm} & \blank{1cm} & \blank{1cm} & \blank{1cm} & \blank{1cm} \\
\hline
$g(x) = f(x+3)$ & \blank{1cm} & \blank{1cm} & \blank{1cm} & \blank{1cm} & \blank{1cm} \\
\hline
\end{tabular}
\end{center}
\smallskip

The output list is \blank{2.6cm} to the parent's. Nothing happened to the
answers at all --- only the \blank{2.4cm} where they appear.

\medskip
\textbf{Where did the vertex go?} The parent's vertex is the point where $f$
was handed $0$. So $g$'s vertex sits at the $x$ that makes the question come
out $0$:

\begin{work}
  x + 3 &= 0 \\
      x &= -3
\end{work}

The graph moved \blank{1.4cm} 3 --- the direction the ``$+$'' does
\textit{not} suggest. This is the harbor again: asking $S$ about 3 a.m.\ means
asking $D$ about 6 a.m., so everything $D$ does, $S$ did three hours
\blank{2cm}.

\smallskip
\begin{center}
\begin{tikzpicture}[x=0.60cm, y=0.34cm]
  \draw[linegray, very thin] (-6,0) grid (3,7);
  \draw[->, thick] (-6.5,0) -- (3.7,0) node[below right, font=\small] {$x$};
  \draw[->, thick] (0,-0.6) -- (0,7.8) node[above, font=\small] {$y$};
  \foreach \x in {-5,-4,-3,-2,-1,1,2}
    \node[below, font=\tiny, fill=white, inner sep=1pt] at (\x,0) {\x};
  \draw[thick, goldacc, dashed, smooth, samples=60, domain=-5.6:-0.4]
    plot (\x, {(\x+3)*(\x+3)});
  \draw[very thick, plum, smooth, samples=60, domain=-2.6:2.6]
    plot (\x, {\x*\x});
  \fill[plum] (0,0) circle (2.2pt);
  \fill[goldacc] (-3,0) circle (2.2pt);
  \draw[->, thick, charcoal] (-0.25,-1.5) -- (-2.75,-1.5);
  \node[below, font=\scriptsize, charcoal] at (-1.5,-1.6) {3 left};
  \node[font=\scriptsize, goldacc, fill=white, inner sep=1pt] at (-4.6,6.4) {$f(x+3)$};
  \node[font=\scriptsize, plum, fill=white, inner sep=1pt] at (1.9,6.4) {$f(x)$};
\end{tikzpicture}
\end{center}

\medskip
\textbf{(b) The rule.} $g(x) = f(x + h)$ is a \textbf{horizontal translation}
by $-h$ units: $h > 0$ moves the graph \blank{1.4cm}, $h < 0$ moves it
\blank{1.4cm}. Do not read the sign in front of the number --- solve for the
input.

\medskip
\textbf{(c) $g(x) = f(2x)$.} To hand $f$ its $2$, you supply $x = $
\blank{1.2cm} \; (you solved $2x = 6$ on the warm-up for the same reason).

\smallskip
\begin{center}
\renewcommand{\arraystretch}{1.4}
\begin{tabular}{|c|c|c|c|c|c|}
\hline
$x$ & $-1$ & $-0.5$ & 0 & $0.5$ & 1 \\
\hline
handed to $f$: $2x$ & \blank{1cm} & \blank{1cm} & \blank{1cm} & \blank{1cm} & \blank{1cm} \\
\hline
$g(x) = f(2x)$ & \blank{1cm} & \blank{1cm} & \blank{1cm} & \blank{1cm} & \blank{1cm} \\
\hline
\end{tabular}
\end{center}
\smallskip

Every landmark arrives at \blank{2cm} the $x$ it used to. That is a
\textbf{horizontal dilation} by a factor of \blank{1.4cm} --- a squeeze, not a
stretch, even though the number is $2$. And $g(x) = f(-x)$ reflects over the
\blank{1.6cm}-axis.

\medskip
The answers were never touched, so the \blank{1.6cm} does not change. The
\blank{2cm} does.
\end{notesbox}

\vspace{0.1in}

\boxguard[16]
\begin{notesbox}{4. Follow the Landmark}
Combining moves needs no order of operations to memorize. Take any point
$(p, q)$ on the parent and answer the two questions separately.

\medskip
\begin{center}
\begin{tcolorbox}[colback=lilac, colframe=plum, boxrule=0.8pt, arc=2mm,
    left=3mm, right=3mm, top=2mm, bottom=2mm, width=0.94\linewidth]
\textbf{For $g(x) = a\,f(x + h) + k$, the parent's point $(p,\,q)$ lands at}
\[
  \bigl(\,p - h,\;\; a q + k\,\bigr).
\]
\textbf{New input:} solve $x + h = p$. \qquad
\textbf{New output:} do to $q$ what the outside says.
\end{tcolorbox}
\end{center}

\medskip
\textbf{Worked example.} $g(x) = -2f(x + 1) - 3$, with $f(x) = x^2$.

\smallskip
\begin{center}
\renewcommand{\arraystretch}{1.5}
\begin{tabular}{|c|c|c|c|}
\hline
\textbf{parent point} $(p,q)$ & \textbf{new input} $p - 1$ & \textbf{new output} $-2q - 3$ & \textbf{lands at} \\
\hline
$(-2,\,4)$ & \blank{1.4cm} & \blank{1.4cm} & \blank{2.4cm} \\
\hline
$(0,\,0)$ & \blank{1.4cm} & \blank{1.4cm} & \blank{2.4cm} \\
\hline
$(2,\,4)$ & \blank{1.4cm} & \blank{1.4cm} & \blank{2.4cm} \\
\hline
\end{tabular}
\end{center}
\smallskip

\begin{itemize}[itemsep=4pt]
  \item In words: left \blank{1cm}, stretched by \blank{1cm}, flipped over the
        \blank{1.6cm}-axis, down \blank{1cm}.
  \item The parent's vertex $(0,0)$ was its lowest point. On $g$ it is the
        \blank{2cm} point, because $a$ is \blank{2cm}.
\end{itemize}
\end{notesbox}

\vspace{0.1in}

\boxguard[14]
\begin{notesbox}{5. The Two Columns Carry Domain and Range}
The domain is a fact about the \blank{2cm} column; the range is a fact about
the \blank{2cm} column. So a transformation moves whichever one it rewrites.

\medskip
Suppose $f$ has domain $[-2,\,4]$ and range $[0,\,9]$, and
$g(x) = -2f(x + 1) - 3$.

\medskip
\textbf{Domain.} $g$ can only be asked about inputs whose question lands
inside $f$'s domain:

\begin{work}
  -2 &\le x + 1 \le 4 \\
  -3 &\le x \le 3
\end{work}

Domain of $g$: \blank{3cm}

\medskip
\textbf{Range.} Every output $q$ of $f$ becomes $-2q - 3$:

\begin{work}
  -2(0) - 3 &= -3 \\
  -2(9) - 3 &= -21
\end{work}

Range of $g$: \blank{3cm} \quad --- note the endpoints traded places, because
multiplying by a \blank{2cm} number reverses order.

\medskip
\textbf{Rule of thumb.} Changes to the \blank{2cm} move the domain. Changes to
the \blank{2cm} move the range.
\end{notesbox}

\vspace{0.1in}

\boxguard[16]
\begin{practicebox}
\begin{enumerate}[itemsep=8pt]
  \item A function $f$ is given by a table. Its peak is at $(2, 4)$.

        \smallskip
        \renewcommand{\arraystretch}{1.4}
        \begin{tabular}{|c|c|c|c|c|c|}
        \hline
        $x$ & 0 & 1 & 2 & 3 & 4 \\
        \hline
        $f(x)$ & 1 & 3 & 4 & 3 & 1 \\
        \hline
        $g(x) = f(x) + 2$ & \blank{1cm} & \blank{1cm} & \blank{1cm} & \blank{1cm} & \blank{1cm} \\
        \hline
        \end{tabular}
        \smallskip

        \begin{enumerate}[label=(\alph*), itemsep=4pt]
          \item Was that a change to the question or the answer?
                \blank{3cm}
          \item Let $h(x) = f(x - 2)$. Then $h(5) = f($ \blank{1.2cm} $) = $
                \blank{1.2cm}.
          \item $h$'s peak is at ( \blank{1.2cm} , \blank{1.2cm} ).
        \end{enumerate}

  \item The lowest point on the graph of $f$ is $(3, -1)$, and the domain of
        $f$ is $[1, 7]$. Let $g(x) = 4f(x - 2) + 5$.

        \begin{enumerate}[label=(\alph*), itemsep=4pt]
          \item Follow the landmark.

                \begin{work}
                  \text{new input} &= 3 + 2 = 5 \\
                  \text{new output} &= 4(-1) + 5 = 1
                \end{work}

                The lowest point of $g$ is ( \blank{1.4cm} , \blank{1.4cm} ).
          \item Domain of $g$: \blank{3cm}
          \item Which of the two moves changed the domain? \blank{4cm}
        \end{enumerate}
\end{enumerate}
\end{practicebox}

\end{document}
